\chapter{Estimation}\label{ch:estimation}

This chapter describes how the parameters of the model in
Chapter~\ref{ch:model} are estimated. The main tool is the expectation
maximisation scheme, whose justification relies on the likelihood of
Equation~\eqref{eq:loglik}.

\section{Latent variables}\label{sec:latent}

Introduce latent indicators $z_i \in \{1,\dots,k\}$ recording the component
that generated observation $x_i$. The complete-data log-likelihood is
\begin{equation}\label{eq:complete}
  \ell_c(\theta) = \sum_{i=1}^{n}\sum_{j=1}^{k} \mathbf{1}\{z_i=j\}\,
  \bigl[\log \pi_j + \log \phi(x_i;\mu_j,\Sigma_j)\bigr].
\end{equation}
Equation~\eqref{eq:complete} is linear in the indicators, which is what makes
the E-step simple.

\lipsum[19-20]

\section{E-step}\label{sec:estep}

The E-step computes the posterior responsibilities
\begin{equation}\label{eq:resp}
  r_{ij} = \frac{\pi_j\,\phi(x_i;\mu_j,\Sigma_j)}
  {\sum_{l=1}^{k}\pi_l\,\phi(x_i;\mu_l,\Sigma_l)} .
\end{equation}
The responsibilities are the expected values of the indicators under the
current parameters, and they are used in the M-step of Section~\ref{sec:mstep}.

\lipsum[21-22]

\section{M-step}\label{sec:mstep}

Maximising the expected complete-data log-likelihood gives the closed-form
updates
\begin{align}
  \pi_j &= \frac{1}{n}\sum_{i=1}^{n} r_{ij}, \label{eq:upd-pi}\\
  \mu_j &= \frac{\sum_i r_{ij}\,x_i}{\sum_i r_{ij}}, \label{eq:upd-mu}\\
  \Sigma_j &= \frac{\sum_i r_{ij}\,(x_i-\mu_j)(x_i-\mu_j)^\top}{\sum_i r_{ij}}.
  \label{eq:upd-sigma}
\end{align}
Each update in Equations~\eqref{eq:upd-pi}--\eqref{eq:upd-sigma} is available
in closed form, so no inner optimisation is required.

\lipsum[23-24]

\section{Estimators}\label{sec:estimators}

For numerical stability the covariance estimates are regularised by adding a
small multiple of the identity, which is the matrix function
of Equation~\eqref{eq:matfun} applied to the eigenvalues. The regularised
estimator is
\begin{equation}\label{eq:reg}
  \widehat{\Sigma}_j^{\,\varepsilon} = \widehat{\Sigma}_j + \varepsilon I_d ,
  \qquad \varepsilon > 0 .
\end{equation}

\lipsum[25-26]

\section{Further remarks}\label{sec:further-estimation}

The remarks in this section restate the main point of the chapter in more
general terms, so that the argument can be read without the surrounding
derivations.

\lipsum[87-98]
